\documentclass[12pt,a4paper,oneside]{article}

\usepackage[spanish]{babel}

\def\mydate{\leavevmode\hbox{\the\year-\twodigits\month-\twodigits\day}}
\def\twodigits#1{\ifnum#1<10 0\fi\the#1}

\usepackage{fancyhdr}
\usepackage{geometry}
\usepackage{graphicx}
\usepackage{wrapfig}
\usepackage{lastpage}
\usepackage[hidelinks]{hyperref}
\usepackage{authblk}
\usepackage{bookmark}

\usepackage[utf8]{inputenc} % Required for inputting international characters
\usepackage[T1]{fontenc} % Output font encoding for international characters
\usepackage{mathpazo} % Palatino font

\makeatletter
\newcommand{\subtitledoc}[1]{\newcommand{\@subtitledoc}{#1}}
\newcommand{\instituto}[1]{\newcommand{\@instituto}{#1}}
\newcommand{\carrera}[1]{\newcommand{\@carrera}{#1}}
\newcommand{\professor}[1]{\newcommand{\@professor}{#1}}
\newcommand{\catedra}[1]{\newcommand{\@catedra}{#1}}
\newcommand{\curso}[1]{\newcommand{\@curso}{#1}}
\newcommand{\legajo}[1]{\newcommand{\@legajo}{#1}}
\newcommand{\footerauthor}[1]{\newcommand{\@footerauthor}{#1}}
\newcommand{\footerlegajo}[1]{\newcommand{\@footerlegajo}{#1}}
\newcommand{\footercatedra}[1]{\newcommand{\@footercatedra}{#1}}

%Configuracion de hoja (margenes y tamaño)
\geometry{a4paper,margin=1in}
\setlength\headheight{28pt}

%Metadata para pdf
%\hypersetup{
%    pdftitle={\@title},
%    pdfsubject={\@catedra\ - \@subtitledoc},
%    pdfauthor={\@author\ - \@legajo\ - \@curso},
%    pdfkeywords={\@title\ \@catedra\ \@subtitledoc\ \@author\ \@legajo\ \@curso}
%}

%formato de encabezado y pie para todas las paginas.
\fancyhead[L]{
  \begin{minipage}[b]{7.5mm}
    \includegraphics[width=7mm]{imagenes/COPYlogo-utn.eps}
  \end{minipage}
  \begin{minipage}[b]{85mm}
    \textbf{Alumnos: }\@footerauthor\\
    \textbf{Legajos: }\@footerlegajo
  \end{minipage}
}
\fancyhead[R]{
  %\textbf{Comisión:} \@curso 
  \textbf{Fecha:} \mydate \\
  \textbf{Cátedra:} \@footercatedra \text{ (}\@curso\text{)}
}
\fancyfoot[C]{} %eliminar antiguo numero de pagina
\fancyfoot[R]{Página \thepage\ de \pageref{LastPage}}
\renewcommand{\headrulewidth}{0.5pt}
\renewcommand{\footrulewidth}{0.5pt}
\pagestyle{fancy}

\addto\captionsspanish{%
  \renewcommand{\contentsname}%
  {CONTENIDO}%
}

\renewcommand{\maketitle}{%
  \newpage
  \thispagestyle{empty}

  \begin{center}

    \textsc{\LARGE \@instituto}\par
    \vspace{0.5cm}
    \textsc{\Large \@carrera}\par
    \vspace{1cm}
    \includegraphics[width=0.30\textwidth]{imagenes/COPYlogo-utn.eps}\par
    \vspace{1cm}

    \textsc{\Large \@catedra}\par
    \vspace{1.5cm}

    \hrule
    \vspace{0.5cm}
    {\huge\bfseries \@title}\par
    \vspace{0.4cm}
    \textsc{\large \@subtitledoc}
    \vspace{0.5cm}
    \hrule \vspace{\baselineskip}


  \end{center}

  \vspace{2cm}

  {\noindent
    \begin{minipage}[t]{.2\textwidth}
      \raggedright
      \textbf{DOCENTES} \par
      \vspace{\baselineskip}
      %\textbf{CÁTEDRA} \par
      \medskip
      \medskip
      \textbf{COMISIÓN} \par
      \medskip
      \medskip
      \textbf{ALUMNOS} \par
    \end{minipage}%
    \begin{minipage}[t]{.05\textwidth}
      \raggedright
      %\textbf{:} \par
      %~\\
      %~\\
      %\textbf{:} \par
      %\textbf{:} \par
      %\medskip
      %\medskip
      %\textbf{:} \par
    \end{minipage}%
    \begin{minipage}[t]{.55\textwidth}
      \raggedright
      \@professor. \par
      \medskip
      %\@catedra. \par
      \medskip
      \@curso \par
      \medskip
      \medskip
      \@author \par
    \end{minipage}%
    \begin{minipage}[t]{.15\textwidth}
      \raggedright
      \vspace{\baselineskip}
      \vspace{\baselineskip}
      \medskip
      \bigskip
      % \medskip
      \medskip
      \medskip
      \@legajo \par
    \end{minipage}
  }
  \vfill
  \begin{center}
    {\Large\textbf{Córdoba, \@date}} \vspace{\baselineskip}
    %\medskip
    %\textbf{Córdoba, Argentina} \par
    \newpage
  \end{center}
}

\makeatother

\usepackage{amsmath}
\usepackage[nottoc,numbib]{tocbibind}
\usepackage{circuitikz}
\usepackage{tikz}
\usepackage{siunitx}
\usepackage{pgfplots}
\pgfplotsset{compat=1.18}
\usepgfplotslibrary{groupplots}
\usepackage{float}

\instituto{Universidad Tecnológica Nacional\\Facultad Regional Córdoba}
\carrera{Ingeniería Electrónica}
\title{Titulo}
\subtitledoc{Subtitulo}
\catedra{Catedra}
\professor{XXXXXXXXXX XXXXXXXX. \par XXXXXXXXXX XXXXXXXX.}
\curso{XRX}
\author{XXXXX XXXXX, XXXXX XXXXX. \par XXXXX XXXXX, XXXXX XXXXX.}
\legajo{XXXXX\par XXXXX}
\footerauthor{XXXXX, XXXXX}
\footerlegajo{XXXXX, XXXXX}
\footercatedra{XXXX}
\date{\today}

\begin{document} %Maximo 10 paginas
\maketitle
\tableofcontents
\newpage
\section{Introducción}
\section{Marco teorico}
\section{Primera Parte}
\subsection{Circuito}
\begin{center}
  \begin{circuitikz}[american]
    \draw (0,0) node[nujt, yscale=-1](ujt){}
    ;
    \draw (ujt.E) to[nos, mirror, invert, l=$L_2$] ++(-2,0)
    to[R] ++(-2,0)
    to[short] ++(0,1) node[vcc]{$V_{E}$}
    ;
    \draw (ujt.B2) to[R] ++(0,2)
    to[nos, l=$L_1$] ++(0,1)
    to[short] ++(0,1) node[vcc]{$V_{CC}$}
    ;
    \draw (ujt.B1) to[short] ++(0,-1) node[ground];
  \end{circuitikz}
\end{center}
\subsection{Procedimiento}
\begin{enumerate}
  \item Armar el circuito seleccionando un correcto valor de las resistencias en
    función del datasheet del UJT.
  \item Abrir el interruptor L1 y cerrar el interruptor L2.
  \item Variar la $V_E$ desde 0-30V y medir la corriente $I_E$.
  \item Completar la tabla propuesta modificándola si fuera necesario.
  \item Graficar la curva $I_E= f(V_E)$ con los datos relevados de la tabla.
  \item Abrir el interruptor L2 y cerrar el interruptor L1.
  \item Variar la VCC desde 0-30V y medir la corriente IB.
  \item Completar la tabla propuesta modificándola si fuera necesario.
\end{enumerate}
\subsection{Simulación}
\subsection{Experimental}
\begin{minipage}{0.2\linewidth}
\begin{tabular}{c|c}
  $V_{E}$ &$I_{E}$  \\
  \hline
  0   &0 \\
  2   &0 \\
  4   &0 \\
  6   &0 \\
  8   &0 \\
  10  &0 \\
  12  &0 \\
  14  &0 \\
  16  &0 \\
  18  &0 \\
  20  &0 \\
  22  &0 \\
  24  &0 \\
  26  &0 \\
  28  &0 \\
  30  &0 
\end{tabular}
\end{minipage}
\begin{minipage}{0.8\linewidth}
\centering
\begin{tikzpicture}
  \begin{axis}[
    width=.9\linewidth,
    xlabel = $V_{E}$,
    ylabel = $I_{E}$,
    grid=major,
    ]
    \addplot[color=red,mark=*] coordinates {
      (0 , 0)
      (2 , 0)
      (4 , 0)
      (6 , 0)
      (8 , 0)
      (10, 0)
      (12, 0)
      (14, 0)
      (16, 0)
      (18, 0)
      (20, 0)
      (22, 0)
      (24, 0)
      (26, 0)
      (28, 0)
      (30, 0)
    };  
  \end{axis} 
\end{tikzpicture}
\end{minipage}

\begin{minipage}{0.2\linewidth}
\begin{tabular}{c|c|c}
  $V_{CC}$ &$V_{B}$ &$I_{B}$  \\
  \hline
  0   &0  &0  \\
  2   &0  &0  \\
  4   &0  &0  \\
  6   &0  &0  \\
  8   &0  &0  \\
  10  &0  &0  \\
  12  &0  &0  \\
  14  &0  &0  \\
  16  &0  &0  \\
  18  &0  &0  \\
  20  &0  &0  \\
  22  &0  &0  \\
  24  &0  &0  \\
  26  &0  &0  \\
  28  &0  &0  \\
  30  &0  &0  \
\end{tabular}
\end{minipage}
\begin{minipage}{0.8\linewidth}
\centering
\begin{tikzpicture}
  \begin{axis}[
    width=.9\linewidth,
    xlabel = $V_{B}$,
    ylabel = $I_{B}$,
    grid=major,
    ]
    \addplot[color=red,mark=*] coordinates {
      (0 , 0)
      (2 , 0)
      (4 , 0)
      (6 , 0)
      (8 , 0)
      (10, 0)
      (12, 0)
      (14, 0)
      (16, 0)
      (18, 0)
      (20, 0)
      (22, 0)
      (24, 0)
      (26, 0)
      (28, 0)
      (30, 0)
    };  
  \end{axis} 
\end{tikzpicture}
\end{minipage}
\begin{center}
 \begin{circuitikz}[american]
   \draw(0,0) to[R, l_=$r_{b2}$]++(0,3) to[short,-o]++(0,0) node[anchor=south](){$B_2$};
   \draw (0,0) to[vR, invert, l=$r_{b1}$]++(0,-3) to[short,-o]++(0,0) node[anchor=north](){$B_1$};
   \draw(-3,0) to[short,-o]++(0,0) node[anchor=east](){$E$} to[D](0,0)
 \end{circuitikz} 
\end{center}

\section{Segunda Parte}
\subsection{Circuito}
\begin{center}
  \begin{circuitikz}
    \draw (0,0) node[nujt, yscale=-1, xscale=-1](ujt){}
    ;
    \draw (ujt.B2) to[R, n=r1, l=100\unit{\ohm}] ++(0,2) coordinate(R100);
    \draw (ujt.B1) to[R, l=22\unit{\ohm}]++(0,-2) 
    to[short] ++(0,-0.5) node[ground]{}
    ;
  \draw (ujt.B2) to[short, -o]++(-1,0) node[anchor=east]{$OUT_1$};
  \draw (ujt.B1) to[short, -o]++(-1,0)node[anchor=east]{$OUT_2$};
  \draw (ujt.E) to[short,-*]++(2,0) coordinate(E) to[C, l=47nF]++(0,-2) node[ground]; 
  \draw (E) to[R, l=4.7K]++(0,2) to[pR, n=POT, l_=100k]++(0,2) coordinate(RPOT) to[short]++(0,1) node[vcc](VCC){$V_{CC}$};
  \draw (POT.wiper) to[short]++(-0.5,0) to[short]++(0,1) to[short, -*](POT|-RPOT);
  \draw (R100) to[short]++(0,1) node[vcc](){$V_{CC}$};
  \end{circuitikz}
\end{center}
\subsection{Procedimiento}
\begin{enumerate}
  \item Armar el circuito.
  \item Medir y graficar la señal en OUT1
  \item Medir y graficar la señal en OUT2
  \item Variar el potenciómetro y observar el efecto sobre la OUT1 y la OUT2 
\end{enumerate}
\subsection{Simulación}
\begin{center}
\begin{tikzpicture}
    \begin{groupplot}[
        group style={
            group size=1 by 4,
            x descriptions at=edge bottom,
            vertical sep=3pt,
        },
        width=0.95\linewidth,
        height=0.4\linewidth,
        grid=both,
        xlabel=$t$,
        thick
    ]
    \nextgroupplot[
      ylabel=$R_{POT}$,
      ]
        \addplot [red]  table [x index=0, y index=4]{./inc/sim/SegundaParte.txt};

    \nextgroupplot[
      ylabel=$V_{C}$,
      ytick={0,2,...,12}
      ]
        \addplot [blue]  table [x index=0, y index=3]{./inc/sim/SegundaParte.txt};
    \nextgroupplot[
      ylabel=$V_{B1}$,
      ytick={0,1,...,4}
      ]
        \addplot [teal]  table [x index=0, y index=1]{./inc/sim/SegundaParte.txt};
    \nextgroupplot[
      ylabel=$V_{B2}$,
      ]
        \addplot [violet]  table [x index=0, y index=2]{./inc/sim/SegundaParte.txt};
    \end{groupplot}
\end{tikzpicture}
\end{center}

\subsection{Experimental}

\section{Tercer Parte}
\begin{center}
  \begin{tabular}{l|c}
    Parametro &Valor \\ 
    \hline
    $\eta$          &  \\
    $R_{BBO}$       &  \\
    $V_{EB1(SAT)}$  &  \\
    $V_{(BR)B1E}$   &  \\
    $P_D$           &  \\
    $I_J$           &  
 \end{tabular} 
\end{center}

\section{Conclusión}
\end{document}
